% RESEARCH_DRAFT_SHA256: cdcacd203d6952447139324986287bb24371cf6eff14baf25b6196ad24b535d4
% Layout and source-reference wrappers expanded for the web renderer.
\documentclass[11pt]{article}
\usepackage[margin=25mm]{geometry}
\usepackage{fontspec,unicode-math}
\setmainfont{Latin Modern Roman}
\setmathfont{Latin Modern Math}
\usepackage{amsmath,amssymb,mathtools,xurl,hyperref}
\hypersetup{colorlinks=true,urlcolor=blue,linkcolor=blue}
\setcounter{secnumdepth}{0}
\setlength{\parindent}{0pt}
\setlength{\parskip}{5pt}
\setlength{\emergencystretch}{3em}
\newcommand{\sref}[2]{\hyperlink{source:#1:#2}{[S#2]}}
\newcommand{\cataloguescope}[1]{\paragraph{Source record.}#1}
\newcommand{\floor}[1]{\left\lfloor#1\right\rfloor}
\begin{document}
\section*{Research-attempt provenance}
This standalone research attempt retains the complete original definition body below,
including its historical source and preparation notes. Those notes describe the
earlier catalogue preparation; the new attempt and its verification are identified
separately after that body. No original ID, catalogue status, or source statement is changed.

\section{Inverse-Polynomial Influence Thresholds on the Boolean Cube}\label{statement}

\subsection{Definitions and mathematical statement}

\paragraph{Shared notation.}Standard number systems are $\mathbb N,\mathbb Z,\mathbb Q,\mathbb R,\mathbb C$. Specialized notation and index ranges follow the cited source and the definitions below.


Write $f(x)=\sum_S\widehat f(S)\prod_{i\in S}x_i$ on the Boolean cube. Degree bounds the largest $|S|$ with nonzero coefficient. A variable\textquotesingle s influence uses the squared Fourier mass on sets containing it, normalized by the source\textquotesingle s total-mass convention.
\paragraph{Record scope and qualifications.}
The catalogue title is retained as metadata. Named certificates, contracts, bounds or frameworks in that title are not asserted as results of the original source extract.
The relative-influence normalization and treatment of zero polynomials are not specified. Without nonzero/normalization assumptions the displayed conclusion has trivial failures; these cases are not silently removed.

Source-record limitation: The relative-influence normalization and nonzero polynomial assumptions are absent; the literal target has unresolved trivial-failure cases.
\paragraph{Statement recorded in the supplied source.}
Let $f$ and $g$ be degree-$d$ multilinear real polynomials on $\{-1,1\}^n$ such that every variable has relative influence at most $\delta$ in both polynomials. Can $\delta$ be chosen inverse-polynomially in $d$ while still guaranteeing a point $x\in\{-1,1\}^n$ for which $f(x)\neq0$ and $g(x)\neq0$? \href{https://doi.org/10.4171/owr/2025/9}{[S2]}
\paragraph{Source status.}
The supplied catalogue labels this source record open; that metadata was checked against the supplied snapshot on 2026-10-11. The label is preserved for provenance and is not an independent verification of present open status or mathematical correctness. The catalogue literature assessment is dated 2026-08-22.
\subsection{Short English statement}
Replace the stated small-influence requirement by an inverse-polynomial degree threshold while guaranteeing simultaneous nonvanishing.
\subsection{Sources}
\begin{itemize}
\item[S1] Ulam AI and the UnsolvedMath Contributors, UnsolvedMath: A Curated Collection of Open Mathematics Problems, catalogue entry. Catalogue curation and metadata: CC BY 4.0. The source extract is retained; definitions, English summaries and qualifications are editorial additions.
\url{https://huggingface.co/datasets/ulamai/UnsolvedMath}.
\item[S2] MFO-RIMS Tandem Workshop: Optimization, Theoretical Computer Science and Algebraic Geometry: Convexity and Beyond (2025), Oberwolfach Reports, DOI 10.4171/owr/2025/9, pp. 405--450
\url{https://doi.org/10.4171/owr/2025/9}.
\item[S3] Austrin--Chung--Chung--Fu--Lin--Mahmoody
\url{https://doi.org/10.1007/978-3-031-15979-4_6}.
\item Source-specific Oberwolfach Reports licences, including CC BY-SA 4.0 where applicable, are retained; see THIRD\_PARTY\_NOTICES.md.
\end{itemize}
\paragraph{Preparation and review.}
The original source and UnsolvedMath contributors retain attribution for the source extract. Codex (GPT-6; exact serving revision not exposed) assisted with definitions, English summaries, scope notes and formatting. No new human mathematical authorship or independently verified solution is asserted. Catalogue identity checked against the complete supplied registry and current repository. Distinct upstream records, including repeated questions and broad or damaged extracts, are retained for complete coverage of the source CS category. A repeated record is not a new independent conjecture. Complete mathematical-equivalence and current-literature review has not been performed. Preparation and final review follow contribution prompts C and D, with source gaps disclosed. Structural and rendering checks do not establish a new theorem or open status.
\label{end:statement}

\section{Literature and context check}
The primary problem session on printed p. 445 defines relative influence as squared coefficient mass containing a variable divided by total squared coefficient mass, including the constant term in the denominator. It reports an exponential sufficient threshold and an inverse-quadratic obstruction, and asks for polynomial dependence on degree. The following degree-one calculation and amplification tests do not improve those general bounds.
\section{Research attempt: The affine threshold and failed product amplification}

Use the uniform Boolean cube and nonzero real multilinear polynomials, so relative influences are defined. For affine functions
\[
f=a+\sum_i b_i x_i,\qquad g=c+\sum_i d_i x_i,
\]
we prove that if they never are simultaneously nonzero, then each has a variable of relative influence $1/2$. Thus a common nonzero point is guaranteed when all influences of both functions are strictly below $1/2$, and this degree-one threshold is sharp.

A nonconstant affine function vanishes on at most half of the cube: fix a coordinate with nonzero coefficient; at most one of its two signs can be a zero for each assignment of the other coordinates. A nonzero constant never vanishes, and would force the other function identically zero if the supports were disjoint. Therefore both functions are nonconstant. Disjoint nonzero supports, each of measure at least $1/2$, force both supports to have measure exactly $1/2$.

Choose $b_i\ne0$. Equality in the preceding count means that for every assignment of the other coordinates,
\[
a+\sum_{j\ne i}b_jx_j\in\{-b_i,b_i\}.
\]
An affine function with at least two nonzero variable coefficients has at least three distinct values on its cube: fixing other coordinates, the two varying terms give extreme values differing by $2(|u|+|v|)$ and a value strictly between the extremes. Hence there is at most one nonzero coefficient besides $b_i$. If there is none, the displayed condition gives $a=\pm b_i$, and the relative influence is $b_i^2/(a^2+b_i^2)=1/2$. If there is one, say $b_j$, the two distinct values $a\pm b_j$ must equal the set $\{-b_i,b_i\}$, so $a=0$ and $|b_j|=|b_i|$; each of the two influences is again $1/2$. Apply the same proof to $g$. Sharpness follows from $f=1+x_1$, $g=1-x_1$, which have disjoint supports and relative influence $1/2$.

\paragraph{Why independent multiplication does not lower influence.}
For nonzero multilinear functions $F$ and $H$ on disjoint variable blocks, the Fourier coefficients of $FH$ are products. Therefore
\[
\sum_S \widehat{FH}(S)^2=\Big(\sum_A\widehat F(A)^2\Big)
\Big(\sum_B\widehat H(B)^2\Big).
\]
For a variable in the $F$ block, restricting the first sum to coefficients containing that variable gives the same second factor, which cancels in the relative influence. Thus multiplying independent copies raises degree but does not reduce any block's maximum relative influence. It cannot turn the sharp affine disjoint-support pair into a low-influence high-degree obstruction.

Averaging independent copies does lower some normalized coordinate weights, but loses disjointness. With $F=\sum_{i=1}^N(1+x_i)$ and $G=\sum_{i=1}^N(1-x_i)$, every summand pair has disjoint supports, yet
\[
FG=N^2-\Big(\sum_i x_i\Big)^2>0
\]
whenever the signs are mixed, for $N\geq2$. Cross terms invalidate the proposed transfer of the no-common-support property. A useful polynomial-degree route must control these cross terms rather than treating them as absent.
\section{Verification and adversarial audit}
The proof uses the total Fourier mass including constants, as specified by the primary report. Equality of support measures is derived before classifying affine zero sets. Strictness is necessary at $1/2$. The product calculation requires disjoint variable blocks; the averaging counterexample is evaluated at mixed signs and includes $N\geq2$.
\paragraph{Audit corrections.} Excluded the zero polynomial, whose relative influence denominator vanishes, and retained the constant Fourier coefficient in every normalization.
\paragraph{External computation.} None; symbolic arguments and boundary checks only.
\section{Scope and remaining gaps}
The inverse-polynomial threshold for arbitrary degree is not established. The sharp affine analysis and two failed amplification operations leave the central degree dependence unchanged.
\paragraph{Integration review.} Codex (GPT-6; exact serving revision not exposed) reviewed the written arguments and expanded layout and source-reference wrappers. Repository metadata and the production web renderer were checked using Python 3.13 and JavaScript with MathJax 3.2.2. This model-assisted review is not independent human proof verification; no PDF compilation or formal proof checking is asserted.

\section{Completion Estimate}
\noindent\textbf{COMPLETION: 0\%}
The calculations below the preserved source statement investigate a concrete route,
but establish no substantiated new progress toward the entire original target.
Known special cases and failed sufficient conditions are not credited as a solution.
\section{References}
The original source references and catalogue attribution are preserved in the
Sources section above. The following URLs were actually consulted for this attempt;
retention of an additional original reference does not claim fresh consultation of it.
\begin{itemize}
\item \url{https://ems.press/content/serial-article-files/51353} (consulted 2026-10-11). 
\end{itemize}
\section{Model and contributor disclosure}
Codex GPT-6 completed written-proof review of this attempt at its declared scope; this is model-assisted review, not independent human proof checking.
OpenAI gpt-6-astra, requested reasoning configuration ultra, generated this research
attempt on 2026-10-11. The exact serving revision is not exposed. Codex GPT-6
performed orchestration and repository-format review; its exact serving revision
is also not exposed. The mathematical route, boundary checks and final adversarial
self-audit followed contribution prompts A then B in
\texttt{docs/contribution\_prompts.md}, after reading \texttt{CONTRIBUTING.md}.
Python 3.14.2 assembled the standalone documents; any mathematical computation is
specified in Verification. No theorem prover was used. No human mathematical
authorship, substantial human proof correction, or independent human proof checking
is asserted. Model-assisted review, source retrieval, and structural checks do not
establish correctness or novelty. Full private reasoning transcripts are unavailable
as reproducibility artifacts; the complete arguments, inputs and outcomes actually
used are stated here or in the companion batch audit.
\section{Final status}
Unresolved as a new research contribution. This is a concrete route investigation,
not a claim of new complete resolution or a current-status certification. A known
answer reproduced for context does not become a new contribution. The catalogue's
source-status label is preserved; this attempt status does not assert that a
source-labelled solved question is open.
\end{document}
